\documentclass[11pt]{article}

% --- TMLR-like single-column formatting ---
% Replace with \usepackage{tmlr} when .sty is available
\usepackage[margin=1.25in]{geometry}
\usepackage{amsmath,amssymb,amsfonts,amsthm}
\usepackage{booktabs}
\usepackage{graphicx}
\usepackage{hyperref}
\usepackage{natbib}
\usepackage{enumitem}
\usepackage{xcolor}
\usepackage{multirow}

\newcommand{\Gr}{\mathrm{Gr}}
\newcommand{\R}{\mathbb{R}}
\newcommand{\misscite}[1]{\textcolor{red}{[CITE: #1]}}

\hypersetup{colorlinks=true, linkcolor=blue!60!black, citecolor=blue!60!black, urlcolor=blue!60!black}

\title{When Does Geometry Help Causal Inference?\\
Boundary Conditions for Sheaf Cohomology, Discrete Curvature,\\
and Subspace Methods in Clinical Epidemiology}

\author{
Elliot Tower
}

\date{July 2026 --- Working Draft (TMLR format)}

\begin{document}
\maketitle

\begin{abstract}
Geometric and topological methods---sheaf cohomology, Grassmannian holonomy, discrete curvature---are increasingly proposed for causal inference and clinical epidemiology, but when they provide genuine advantages over standard statistical tools remains unclear.
We characterize this boundary through eleven simulation experiments spanning federated data consistency, confound detection, causal graph validation, and treatment heterogeneity, with multiple sclerosis neurology as a motivating application.

We identify three structural conditions under which geometry adds signal.
First, \emph{subspace-valued data}: Grassmannian holonomy detects global inconsistency invisible to all pairwise and scalar alternatives (Berry phase $= 1.85$, $p < 0.001$, while maximum pairwise distance remains below the noise threshold), but sheaf consistency testing on scalar data reduces algebraically to Cochran's $Q$.
Second, \emph{cyclic consistency constraints}: composed parallel transport accumulates signal as $\sqrt{m}$ over $m$ steps while noise grows as a random walk, a separation mechanism unavailable to pairwise tests.
Third, \emph{edge-specific heterogeneity}: per-edge sheaf $Q$ tests recover DAG structure that global tests miss ($p < 10^{-300}$ vs.\ $p = 0.659$), and $H^1$ effect-modifier classification achieves perfect transportability prediction with a three-order-of-magnitude gap between categories.

When these conditions are absent, geometry reduces to standard tools or fails outright.
We also document two failure modes---Ollivier--Ricci curvature fails for edge validation (AUROC $= 0.466$) while Forman--Ricci succeeds (AUROC $= 0.677$), and PCA-based dimensionality reduction destroys treatment effect signal regardless of CATE estimator quality---providing practical guidance for method selection.
\end{abstract}

%% ==============================================================
\section{Introduction}
\label{sec:intro}
%% ==============================================================

Multiple sclerosis (MS) illustrates structural challenges common across clinical epidemiology.
Disease mechanisms vary across patient strata---relapsing vs.\ progressive phenotypes, treatment regimens, genetic backgrounds---yet some causal relationships remain consistent across all strata, such as the downstream path from neurodegeneration to disability \misscite{Lublin 2014 phenotype classification}.
Federated imaging studies across hospital sites carry acquisition-dependent confounds \misscite{multi-site MRI harmonization}.
Causal DAGs estimated from observational data contain both true edges and spurious associations with no reliable post-hoc filter \misscite{causal discovery review}.
Treatment heterogeneity---the fact that patients respond differently to the same intervention---remains difficult to detect and characterize \misscite{Athey \& Imbens HTE review}.

Each problem has a geometric aspect.
Federated consistency asks whether local estimates are compatible under known transformations---what sheaf cohomology measures \misscite{Robinson 2014; Curry 2014 sheaves}.
For multivariate data, the relevant geometry lives on the Grassmannian $\Gr(k,d)$, whose curvature generates Berry phase when local sections are transported around closed loops \misscite{Berry 1984; Simon 1983}.
Edge validation in causal graphs asks whether topological features carry information about edge truthfulness.
Treatment subtype recovery asks whether patients cluster in CATE-augmented covariate space.
The question is whether these connections produce practical advantages or repackage standard statistics in topological language.

We answer this by characterizing the boundary between the two regimes.
The paper makes three contributions:
\begin{enumerate}[nosep]
    \item We identify three structural conditions---subspace-valued data, cyclic consistency constraints, edge-specific heterogeneity---under which geometric methods detect structure invisible to scalar and pairwise alternatives.
    \item We show that absent these conditions, geometric methods reduce to standard tools (sheaf tests $=$ Cochran's $Q$, partial correlation $=$ perfect confound detection) or fail outright (ORC curvature, PCA-based HTE clustering).
    \item We validate the positive cases on MS-calibrated simulations: Grassmannian holonomy, per-edge sheaf DAG testing, and $H^1$ effect-modifier classification each detect heterogeneity structure that existing methods miss.
\end{enumerate}

%% ==============================================================
\section{Methods}
\label{sec:methods}
%% ==============================================================

\subsection{Sheaf consistency testing}
\label{sec:meth-sheaf}

A sheaf on a network $G = (V, E)$ assigns local data (\emph{stalks}) to each vertex and requires that overlapping assignments agree on shared edges via \emph{restriction maps}.
When they disagree, the obstruction lives in the first cohomology group $H^1$, which measures the failure of global consistency \misscite{Curry 2014; Robinson 2014}.

\paragraph{Scalar stalks.}
For scalar stalks (effect estimates $\hat\beta_k$ with standard errors $\mathrm{se}_k$) and pairwise-difference restriction maps on a complete graph of $K$ sites, the sheaf Laplacian test statistic is
\begin{equation}
    Q = \mathbf{o}^\top \Sigma^{-1} \mathbf{o}, \quad
    \mathbf{o} = D_0 \boldsymbol{\hat\beta}, \quad
    \Sigma = D_0 \operatorname{diag}(\mathrm{se}^2) D_0^\top,
    \label{eq:sheaf-Q}
\end{equation}
where $D_0$ is the signed incidence matrix and $Q \sim \chi^2_{\mathrm{rank}(D_0)}$ under the null.

\paragraph{Subspace-valued stalks.}
For sections $V_1, \ldots, V_m \in \Gr(k,d)$ arranged in a cycle, the parallel transport from $V_i$ to $V_{i+1}$ is $T_{i \to i+1} = U V^\top$ where $V_i^\top V_{i+1} = U \Sigma V^\top$ (SVD).
The composed \emph{holonomy} $\Phi = T_{m \to 1} \cdots T_{1 \to 2}$ equals the identity if and only if sections are globally consistent.
For the linked-column construction (two columns of $V_0$ rotating toward shared perpendicular directions with $\pi/2$ phase offset at radius $r$), the holonomy norm is
\begin{equation}\label{eq:berry}
    \|\Phi - I_k\|_F = 2\sqrt{2}\,|\sin(\pi \sin^2 r)|.
\end{equation}

\paragraph{Per-edge testing.}
For a causal DAG estimated in $S$ strata, each edge is tested independently:
$Q^{(e)} = \mathbf{z}^\top \Sigma^{-1} \mathbf{z}$ with $z_{ij} = \hat\beta_i^{(e)} - \hat\beta_j^{(e)}$, under $Q^{(e)} \sim \chi^2_{S-1}$.
Bonferroni correction controls the family-wise error rate.

\paragraph{$H^1$ effect-modifier classification.}
For a mechanism--modifier pair, patients are stratified by the modifier and the causal effect estimated within each stratum.
The sheaf $Q$ test classifies the pair as transportable ($Q$ non-significant, $H^1 \approx 0$) or stratum-specific ($Q$ significant, $H^1 \neq 0$).
Sample size per stratum is calibrated to expected interaction strength.

\subsection{Confound detection}
\label{sec:meth-confound}

\paragraph{Partial correlation collapse.}
For each biomarker $j$, we compute raw Pearson correlation, partial correlation (conditioning on all other covariates), and intensive margin correlation (conditioning on treatment).
Confounded associations should collapse under partial conditioning.

\paragraph{Bracket-norm confound audit.}
For multi-site imaging, confound leakage is $\Delta = (R^2_{\mathrm{metric}} - R^2_{\mathrm{unique}}) / R^2_{\mathrm{metric}}$, where $R^2_{\mathrm{unique}} = R^2_{\mathrm{both}} - R^2_{\mathrm{acq}}$ is disease-severity variance unexplained by acquisition parameters.

\subsection{Curvature-based edge validation}
\label{sec:meth-curvature}

We test whether edge features discriminate true positive (TP) from false positive (FP) edges in graphs learned from structural equation models.
Six features are tested: Forman--Ricci curvature ($\kappa_F = 4 - d(u) - d(v)$) \misscite{Forman 2003}, augmented Forman ($\kappa_F + 3|\mathcal{N}(u) \cap \mathcal{N}(v)|$), Jaccard coefficient, edge betweenness \misscite{Girvan \& Newman 2002}, partial correlation magnitude, and average endpoint clustering coefficient.
Ollivier--Ricci curvature \misscite{Ollivier 2009; Ni et al.\ 2019} serves as a baseline.
Discrimination is measured by AUROC across 6 conditions (linear/nonlinear SEM $\times$ 3 thresholds, 60 graphs per condition).

\subsection{Treatment heterogeneity detection}
\label{sec:meth-hte}

Four CATE estimators (KNN, random forest, gradient boosting T-learners; random forest S-learner) are crossed with four clustering methods (PCA $+$ $K$-means, CATE $K$-means, CATE$+$covariates $K$-means, spectral on CATE RBF kernel).
Oracle baselines (true CATE) and a naive baseline (raw-covariate $K$-means) bound performance.

%% ==============================================================
\section{Results}
\label{sec:results}
%% ==============================================================

Results are organized by the three boundary conditions rather than by experiment batch.
For each condition, we show the case where geometry reduces to standard tools and the case where it adds signal.

\subsection{Condition 1: Scalar vs.\ subspace-valued data}
\label{sec:res-scalar-vs-subspace}

\paragraph{Scalar stalks reduce to Cochran's $Q$.}
We simulate $K = 8$ hospital sites with scalar treatment effect estimates on a complete graph.
The sheaf test achieves proper calibration (type~I error $= 5.85\%$ at $\alpha = 0.05$) with power reaching $96.8\%$ at bias $\delta = 0.20$, and localizes the biased site to rank~1 in 94\% of simulations.
At every bias level, structured separation, and replicate, the sheaf test produces results \emph{identical} to Cochran's $Q$.
On a complete graph with scalar stalks and pairwise-difference restriction maps, the quadratic form in Eq.~\eqref{eq:sheaf-Q} is algebraically equivalent to $Q$.

\paragraph{Subspace-valued stalks produce genuine holonomy.}
We test on $\Gr(3, 20)$ with $m = 24$ sections forming a closed loop, using a linked-column Berry phase construction at radius $r = 0.5$.
The cocycle obstruction test evaluates three simultaneous conditions: (C1)~pairwise consistency, (C2)~global inconsistency, (C3)~scalar blindness.

\begin{table}[t]
\centering
\caption{Cocycle obstruction: pairwise consistent yet globally inconsistent.}
\label{tab:cocycle}
\small
\begin{tabular}{@{}lcc@{}}
\toprule
\textbf{Metric} & \textbf{Value} & \textbf{Threshold} \\
\midrule
Measured holonomy $\|\Phi - I\|_F$ & $1.851$ & --- \\
Predicted (Eq.~\ref{eq:berry}) & $1.869$ & --- \\
$p$-value vs.\ null & $< 0.001$ & --- \\
\midrule
Max pairwise distance & $0.730$ & $0.767$ (C1 holds) \\
Holonomy norm & $1.851$ & $1.226$ (C2 holds) \\
\bottomrule
\end{tabular}
\end{table}

All three conditions hold simultaneously (Table~\ref{tab:cocycle}).
The measured holonomy ($1.851$) matches the predicted Berry phase ($1.869$) to within 1\%.
Per-step rotation is masked by noise (max pairwise distance $0.730 < 0.767$), but the composed holonomy accumulates coherently far above the null threshold ($1.851 > 1.226$).

Three competitor baselines fail.
Maximum pairwise angle cannot distinguish planted from control sections ($0.730$ vs.\ $0.695$, 5\% difference).
Random-effects and CKA tests detect subspace spread ($p < 10^{-77}$) but are blind to the cyclic obstruction that makes the holonomy non-trivial.
Only composed parallel transport detects the global inconsistency.

\textbf{Boundary.} The transition from scalar to subspace-valued data is the primary determinant of whether geometric structure carries information.
The holonomy signal accumulates coherently as $\sqrt{m}$ over $m$ loop steps, producing a signal-to-noise ratio of $\sqrt{m} \approx 5$ for $m = 24$---a separation mechanism with no scalar analogue.

\subsection{Condition 2: Global vs.\ edge-specific heterogeneity}
\label{sec:res-global-vs-edge}

\paragraph{Global tests miss edge-specific structure.}
A 3-node MS DAG (inflammation $\leftrightarrow$ degeneration $\to$ disability) is estimated in 8 disease strata.
The feedback edges vary dramatically: early RRMS shows dominant infl$\to$degen ($0.404$) with negligible reverse flow ($-0.001$), while PPMS shows the opposite ($-0.008$ and $0.385$).
Disability edges remain stable (variance $\sim\!100\times$ smaller).

A global Frobenius norm test averages these variances and produces $p = 0.659$ (non-significant).
Per-edge sheaf $Q$ tests recover the true structure (Table~\ref{tab:dag}).

\begin{table}[t]
\centering
\caption{Per-edge sheaf $Q$ tests on the MS inflammation--degeneration--disability DAG across 8 strata.}
\label{tab:dag}
\small
\begin{tabular}{@{}lcccl@{}}
\toprule
\textbf{Edge} & $\boldsymbol{Q}$ & $\boldsymbol{p}$ & \textbf{df} & \textbf{Het.?} \\
\midrule
infl $\to$ degen & $1806.9$ & $< 10^{-300}$ & 7 & Yes \\
degen $\to$ infl & $1549.6$ & $< 10^{-300}$ & 7 & Yes \\
infl $\to$ disab & $7.05$ & $0.423$ & 7 & No \\
degen $\to$ disab & $9.98$ & $0.189$ & 7 & No \\
\bottomrule
\end{tabular}
\end{table}

Treatment signatures align with known mechanisms: BTK inhibition suppresses infl$\to$degen from $0.404$ to $0.002$; siponimod preserves degen$\to$disab ($0.386$), consistent with relapse-independent progression \misscite{Kappos 2018 siponimod}.

\paragraph{$H^1$ classification achieves clean bimodal separation.}
Seven MS mechanism--modifier pairs are classified as transportable or stratum-specific (Table~\ref{tab:h1}).

\begin{table}[t]
\centering
\caption{$H^1$ effect-modifier classification: three-order-of-magnitude gap between categories.}
\label{tab:h1}
\small
\begin{tabular}{@{}llccrl@{}}
\toprule
\textbf{Pair} & \textbf{Expected} & $\boldsymbol{\gamma}$ & $\boldsymbol{Q}$ & $\boldsymbol{p}$ & \textbf{Correct} \\
\midrule
HLA $\times$ EBV & transport & $0.10$ & $3.30$ & $0.348$ & Yes \\
EBV necessity & transport & $0.05$ & $6.47$ & $0.091$ & Yes \\
Vitamin D & transport & $0.04$ & $6.97$ & $0.073$ & Yes \\
\midrule
Sex $\times$ course & non-transp & $0.50$ & $2434$ & $< 10^{-300}$ & Yes \\
Genetics $\times$ OCB & non-transp & $0.60$ & $3588$ & $< 10^{-300}$ & Yes \\
Age $\times$ anti-CD20 & non-transp & $0.40$ & $1705$ & $< 10^{-300}$ & Yes \\
Phenotype $\times$ GM & non-transp & $0.45$ & $1828$ & $< 10^{-300}$ & Yes \\
\bottomrule
\end{tabular}
\end{table}

All 7 pairs correctly classified (100\%).
Transport pairs ($Q < 7$, $p > 0.07$) are exposure$\to$mechanism edges with weak interaction ($\gamma \leq 0.10$); non-transport pairs ($Q > 1700$, $p < 10^{-300}$) have strong interactions ($\gamma \geq 0.40$).

\textbf{Boundary.} When a DAG contains both heterogeneous and homogeneous edges, global tests dilute the signal.
Per-edge sheaf $Q$ tests avoid this by testing each structural relationship independently.
The $H^1$ classifier operationalizes this as a practical workflow: classify each mechanism--modifier pair as transportable or stratum-specific with no manual threshold tuning.

\subsection{Condition 3: Method-specific failure modes}
\label{sec:res-failures}

Two negative results characterize failure modes orthogonal to the boundary conditions above.

\paragraph{Curvature type: ORC fails, Forman--Ricci succeeds.}
Ollivier--Ricci curvature yields AUROC $= 0.466$ for TP vs.\ FP edge discrimination (below chance).
Forman--Ricci curvature rescues the approach at AUROC $= 0.677$ (Table~\ref{tab:curvature}), consistent across all six conditions ($0.634$--$0.677$).
Partial correlation magnitude provides a complementary non-topological signal (AUROC $= 0.657$).

\begin{table}[t]
\centering
\caption{Edge feature AUROC for TP vs.\ FP discrimination in learned causal graphs (best condition per feature, 60 graphs/condition).}
\label{tab:curvature}
\small
\begin{tabular}{@{}lccc@{}}
\toprule
\textbf{Feature} & \textbf{Best AUROC} & \textbf{Condition} & \textbf{Direction} \\
\midrule
Forman--Ricci & $\mathbf{0.677}$ & linear / loose & TP $>$ FP \\
$|\text{Partial corr}|$ & $0.657$ & nonlinear / loose & TP $>$ FP \\
Betweenness & $0.610$ & linear / loose & TP $>$ FP \\
Avg.\ clustering & $0.568$ & nonlinear / loose & TP $>$ FP \\
\midrule
ORC (initial) & $0.466$ & linear & wrong dir. \\
Augmented Forman & $0.484$ & --- & wrong dir. \\
Jaccard & $0.408$ & --- & wrong dir. \\
\bottomrule
\end{tabular}
\end{table}

The mechanism is a degree asymmetry: $\kappa_F = 4 - d(u) - d(v)$ is higher when both endpoints have low degree.
TP edges connect lower-degree nodes on average because real causal relationships produce specific partial correlations, while FP edges arise from indirect paths through higher-degree hub nodes.
ORC measures optimal transport between neighborhoods (global topology), which carries no discriminative information in these sparse graphs.

\paragraph{Dimensionality reduction: PCA destroys treatment effect signal.}
We simulate $n = 3{,}000$ patients with 3 treatment response subtypes (CATE $\in \{+2.0, -2.5, 0\}$).
The $4 \times 4$ CATE $\times$ clustering grid (Table~\ref{tab:hte}) shows that PCA clustering fails universally: all four CATE methods and the oracle (ARI $= 0.095$) produce near-random clusters.

\begin{table}[t]
\centering
\caption{ARI for treatment subtype recovery ($4 \times 4$ grid, 100 replicates, $n = 3{,}000$).}
\label{tab:hte}
\small
\begin{tabular}{@{}lcccc@{}}
\toprule
 & \textbf{PCA} & \textbf{CATE} & \textbf{CATE+cov} & \textbf{Spectral} \\
\midrule
KNN       & $-0.011$ & $0.179$ & $0.225$ & $0.085$ \\
RF T-learn  & $-0.013$ & $0.238$ & $\mathbf{0.270}$ & $0.097$ \\
GBM T-learn & $-0.016$ & $0.189$ & $0.264$ & $0.055$ \\
RF S-learn  & $-0.014$ & $0.175$ & $0.245$ & $0.066$ \\
\midrule
Oracle        & $0.095$  & $1.000$ & $0.551$ & $1.000$ \\
Naive (raw cov) & \multicolumn{4}{c}{$0.218$} \\
\bottomrule
\end{tabular}
\end{table}

The best estimated combination (RF T-learner $+$ covariate-augmented $K$-means, ARI $= 0.270$) exceeds the naive baseline ($0.218$) by 24\% but falls far short of the oracle ($1.000$).
The bottleneck is CATE estimation noise at $n = 3{,}000$: the oracle achieves perfect recovery when true CATE is known.

\subsection{Confound detection}
\label{sec:res-confound}

\paragraph{Partial correlation under complete confounding.}
With 10 real and 10 confounded biomarkers ($n = 3{,}000$), partial correlation achieves AUROC $= 1.000$ at all sample sizes ($n = 200$--$5{,}000$) with zero variance.
Confounded biomarkers collapse to partial correlations of $0.001$--$0.033$; real biomarkers retain $0.61$--$0.76$.
This perfect separation is a property of complete confounding (shared-cause pathways that conditioning fully removes) and would degrade under partial confounding or unobserved confounders.

\paragraph{Bracket-norm audit under acquisition confounds.}
Four MS imaging biomarkers (iron rim QSM, deep GM atrophy, cortical lesion count, cervical cord CSA; $n = 1{,}500$, 8 sites) all pass: confound leakage $\Delta < 0.05$ for three metrics and $\Delta = -0.19$ (suppressor effect) for iron rim QSM.
Post-correction partial correlations with sNfL remain significant ($p < 10^{-55}$).

\subsection{MS prerequisites}

Cohort contamination rate is $0.9\%$ (below 2\% threshold), effect estimates are stable across diagnostic stringency ($p = 0.18$), and IRT-derived theta scores correlate more strongly with biological biomarkers than raw EDSS ($r = 0.836$ vs.\ $0.801$ for sNfL).

%% ==============================================================
\section{Discussion}
\label{sec:discussion}
%% ==============================================================

\subsection{The boundary}

The three conditions we identify---subspace-valued data, cyclic constraints, edge-specific heterogeneity---are sufficient conditions, not a closed characterization.
They share a common structure: each creates information that standard scalar/pairwise/global methods cannot access.
Holonomy accumulates coherently around cycles ($\sqrt{m}$ signal growth) in a way that pairwise comparisons cannot replicate.
Per-edge tests avoid the variance dilution that plagues global tests on heterogeneous DAGs.
Subspace-valued data carry covariance structure that scalar projections discard.

The conditions are also \emph{testable in advance}: a practitioner can determine before running geometric analyses whether their data are multivariate, whether consistency constraints span cycles, and whether heterogeneity is plausibly edge-specific.
When none of these conditions hold, standard tools (Cochran's $Q$, partial correlation, $K$-means) are sufficient and simpler.

\subsection{Curvature selection}

The ORC/Forman--Ricci contrast illustrates a broader principle: the ``right'' curvature depends on which structural property discriminates the classes of interest.
ORC measures optimal transport between neighborhoods, which depends on global topology.
Forman--Ricci measures degree deficit, which captures the local connectivity asymmetry between true and spurious edges.
For edge validation in learned causal graphs, local structure matters and global structure does not.

\subsection{Implications for MS and neuro-epidemiology}

The MS results illustrate how these boundary conditions arise naturally in clinical settings.
The two-process model (inflammation $\leftrightarrow$ degeneration $\to$ disability) generates heterogeneity structure invisible to global tests but recoverable by per-edge sheaf analysis.
Multi-site imaging cohorts produce subspace-valued data (biomarker covariance matrices) where holonomy could detect acquisition-dependent distortions that scalar harmonization methods miss.
Transportability assessment---which effects generalize across patient strata---should be a standard step before pooling multi-site estimates.

These applications are not MS-specific.
Any setting where causal relationships vary across population strata and multivariate biomarker data are available can use the same tools: Alzheimer's disease (ATN framework), oncology (molecular subtype heterogeneity), and psychiatric genetics (gene $\times$ environment interactions).

\subsection{Limitations}

All experiments use simulated data.
The simulations test specific methodological claims rather than modeling realistic clinical complexity: confounding is complete (all-or-nothing), SEMs use random DAGs without domain structure, HTE subtypes are well-separated, and MS heterogeneity patterns are planted.
Real clinical data present partial confounding, measurement error, missing data, and continuous rather than discrete variation.

The cocycle dose-response shows substantial noise at small radii, suggesting that larger $m$ or repeated trials would be needed for reliable dose-response estimation.
The $H^1$ classifier's scale-invariance stress test did not fully pass: quantile normalization reduced scale variability from $0.517$ to $0.346$ but did not reach the $0.30$ target.
Developing scale-free alternatives to OLS-based sheaf $Q$ tests is an open direction.

%% ==============================================================
\section{Conclusion}
\label{sec:conclusion}
%% ==============================================================

Geometric and sheaf-cohomological methods help causal inference under three identifiable conditions: subspace-valued data, cyclic consistency constraints, and edge-specific heterogeneity.
Absent these conditions, they reduce to Cochran's $Q$, partial correlation, or fail at tasks where simpler alternatives succeed.
The boundary itself is the contribution: it tells practitioners when to reach for geometric tools and when standard statistics suffice.

\paragraph{Reproducibility.}
All simulation code and results are available at \textcolor{red}{[REPO URL]}. \textcolor{red}{[Zenodo DOI block to be added at submission.]}

\bibliographystyle{unsrtnat}
% \bibliography{references}

\end{document}
